\section{Setting, requirements and threat model}\label{sec:setting}

\subsection{Actors and roles}

A \emph{person} enrolls once and receives an anonymous credential. From it the person derives
exactly one pseudonym for each of three roles:
\begin{itemize}
  \item \textbf{author}: proposes items;
  \item \textbf{evaluator}: reviews items proposed by others;
  \item \textbf{respondent}: answers trial items, which are mixed with validated ones and do not count
    toward the respondent's own score.
\end{itemize}
The three pseudonyms of one person are unlinkable to each other and to the person. A
\emph{consortium} of independent, always-on signer nodes stores the append-only record of items
and reviews, recomputes the scores and signs the result. Anyone may run a light node that
verifies signatures. A threshold \emph{issuing committee}, distinct from both the consortium and
the state identity provider, issues credentials. The state authenticates a person but never
issues the credential, so it cannot link a person to their activity.

\subsection{Requirements}

The following requirements are the invariants of the design~\cite{isegoria-repo}. Every
mechanism below is evaluated against them.

\begin{requirement}[Anonymity]\label{req:anon}
No personal, demographic or affiliation data enters the system, not even optionally.
\end{requirement}
\begin{requirement}[No majority rule]\label{req:nomajority}
Acceptance of an item is never decided by counting votes.
\end{requirement}
\begin{requirement}[No money]\label{req:nomoney}
There is no token, stake or payment. Proposing costs reputation and is rate-limited.
\end{requirement}
\begin{requirement}[Separated reputations]\label{req:separate}
The author score and the evaluator score live on unlinkable pseudonyms and are never combined.
\end{requirement}
\begin{requirement}[Non-rotatable pseudonyms]\label{req:norotate}
A person can derive exactly one pseudonym per role, so a bad reputation cannot be discarded~\cite{friedman2001pseudonyms}.
\end{requirement}
\begin{requirement}[Reproducibility]\label{req:repro}
The scores are a deterministic function of the recorded inputs, identical bit for bit on every
recomputation.
\end{requirement}
\begin{requirement}[Batch validation]\label{req:batch}
Empirical validation is performed on groups of items, never on a single item.
\end{requirement}

\subsection{Adversaries}

\Cref{tab:threats} lists the adversaries the design addresses and where each is analysed.
The adversary that motivates the whole design is the \emph{majority faction}: a large, cohesive
group of real persons that votes sincerely but one-sidedly. It needs no fraud to capture a
voting system. A \emph{cartel} is a coordinated bloc of real persons. A \emph{Sybil}
adversary creates fake identities. A \emph{long-con} reviewer builds reputation in order to spend
it later. A \emph{biased author} writes items whose wording favours one group.

\begin{table}[t]
\centering\small
\caption{Adversaries, primary countermeasures, and where this paper analyses them.}
\label{tab:threats}
\begin{tabular}{@{}>{\raggedright\arraybackslash}p{0.25\linewidth}>{\raggedright\arraybackslash}p{0.43\linewidth}>{\raggedright\arraybackslash}p{0.24\linewidth}@{}}
\toprule
Adversary & Primary countermeasure & Analysis \\
\midrule
Majority faction & Bridging aggregation (Level A) & \cref{sec:levela} \\
Biased item passing review & Empirical DIF on latent classes (Level B) & \cref{sec:levelb} \\
Cartel & Random assignment; sublinear discount & \cref{sec:collusion} \\
Sybil, whitewashing & External uniqueness; non-rotatable pseudonyms & \cref{ass:unique} \\
Brigading, herding & Unpredictable assignment; blind commit--reveal & \cref{ass:assign,ass:blind} \\
Long-con reviewer & Evaluator score on honeypots; asymmetric updates; cap & \cref{sec:levelc} \\
Dishonest signer & Deterministic recomputation by a threshold consortium & \cref{ass:repro} \\
Sample poisoning & Distinct respondents; batch validation & \cref{sec:limitations} \\
Statistical deanonymization & Out of scope here (text normalization, delays) & \cref{sec:limitations} \\
\bottomrule
\end{tabular}
\end{table}

\subsection{Assumptions that abstract the identity and network layers}\label{sec:assumptions}

The scoring mechanism is the subject of this paper. The cryptographic identity layer and the
replicated storage layer enter only through the following assumptions, so the analysis does
not depend on how they are realized.

\begin{assumption}[Uniqueness and unlinkability]\label{ass:unique}
Each person holds at most one credential. Each credential yields exactly one pseudonym per role.
Pseudonyms of different roles are computationally unlinkable to each other and to the
person's identity.
\end{assumption}

\begin{assumption}[Unpredictable assignment]\label{ass:assign}
Admission, reviewer assignment, honeypot placement and sortition use public randomness that no
participant can predict or bias before the inputs it applies to are fixed.
\end{assumption}

\begin{assumption}[Blind simultaneous review]\label{ass:blind}
Every reviewer of an item commits to a judgment before any judgment is revealed. Reviewers
see neither the author nor the other judgments.
\end{assumption}

\begin{assumption}[Reproducible computation]\label{ass:repro}
The scores are recomputed from the recorded inputs by a threshold of independent signers, and
a result is accepted only if enough signers agree.
\end{assumption}

\begin{assumption}[Respondent sampling]\label{ass:sample}
The respondents of a pilot are distinct persons. Their abilities and latent group memberships
represent the population for which the item bank is intended.
\end{assumption}

In the reference implementation, \cref{ass:unique} is realized with a threshold
oblivious pseudorandom function over the state identity~\cite{rfc9497},
threshold BBS+ blind credentials~\cite{boneh2004short,au2006bbsplus,tessaro2023bbs} and
per-role nullifiers. Two parts are still missing: binding the OPRF input to the
state-authenticated identity (roadmap item T20), and a distributed key generation that removes the
trusted dealer from both committees (T19). \Cref{ass:assign} uses randomness derived from the latest signed
checkpoint; the repository records that this derivation can still be ground by whoever orders
the last deposits (roadmap item T37), so \cref{ass:assign} is not yet met. \Cref{ass:blind}
uses commitments that bind the committer and the item. \Cref{ass:repro} relies on a deterministic
Rust engine (\cref{sec:implementation}). \Cref{ass:sample} is the strongest assumption and the
least under the system's control, because respondents are self-selected members of the network
(\cref{sec:limitations}).

\subsection{Lifecycle of an item}

\Cref{fig:pipeline} shows the lifecycle.
\begin{enumerate}
  \item \textbf{Deposit.} An author deposits a draft item with a mandatory primary source.
  \item \textbf{Admission.} Each epoch a lottery admits a subset of the deposits, because
    validation capacity, not proposals, is the bottleneck.
  \item \textbf{Review.} Each admitted item is reviewed blind by $k=7$--$11$ reviewers,
    assigned at random and stratified on their estimated latent position. Each reviewer
    declares a probability that the item will pass empirical validation; that number is both
    the rating used by Level~A and the forecast scored by Level~C.
  \item \textbf{Bridging gate.} Level~A computes a bridge score $B_j$. The item passes if
    $B_j\ge\tau+\varepsilon$. If $B_j$ falls in the band $[\tau-\varepsilon,\tau+\varepsilon)$,
    the item gets more reviewers and is re-decided against $\tau$. Below the band, it is
    rejected, or, if it was rejected for polarization (large $|f_j|$), its author may appeal
    to the evidence filter.
  \item \textbf{Pilots.} Stage~1 ($\approx300$ respondents) screens out broken and
    non-discriminating items. Stage~2 (batches, $\approx3000$ respondents for latent DIF)
    fits IRT and tests for DIF.
  \item \textbf{Pool.} Surviving items enter the active pool. The pool is periodically
    re-validated as a whole, and items are retired for exposure, drift or emerging DIF.
\end{enumerate}

\begin{figure}[t]
\centering
\begin{tikzpicture}[
  node distance=4mm and 5mm,
  box/.style={draw, rounded corners=2pt, align=center, font=\small, minimum height=8mm, inner sep=3pt, fill=gray!6},
  dec/.style={draw, diamond, aspect=2.2, align=center, font=\small, inner sep=1pt, fill=blue!5},
  arr/.style={-{Stealth[length=2mm]}, thick}]
\node[box] (dep) {Deposit\\(item + source)};
\node[box, right=of dep] (lot) {Admission\\lottery};
\node[box, right=of lot] (rev) {Blind review\\$k$ random reviewers};
\node[dec, right=of rev] (gate) {Bridging\\gate};
\node[box, below=9mm of gate] (p1) {Pilot 1\\screen};
\node[box, left=of p1] (p2) {Pilot 2 (batch)\\IRT + latent DIF};
\node[box, left=of p2] (pool) {Active pool\\re-validation};
\node[box, left=of pool] (ret) {Retirement};
\node[box, right=8mm of gate, fill=red!6] (rej) {Rejected};
\node[box, above=6mm of gate, fill=orange!8] (app) {Appeal to\\evidence};
\node[box, above=6mm of rev, fill=yellow!10] (band) {Supplementary\\review};
\draw[arr] (dep) -- (lot);
\draw[arr] (lot) -- (rev);
\draw[arr] (rev) -- (gate);
\draw[arr] (gate) -- node[right, font=\scriptsize]{pass} (p1);
\draw[arr] (gate) -- node[above, font=\scriptsize]{defect} (rej);
\draw[arr] (gate) -- node[right, font=\scriptsize]{polarized} (app);
\draw[arr] (gate.north west) -- node[above, font=\scriptsize, pos=0.55, sloped]{band} (band.east);
\draw[arr] (band.south) -- (rev.north);
\draw[arr] (app.east) -| ($(rej.east)+(4mm,0)$) |- (p1.east);
\draw[arr] (p1) -- (p2);
\draw[arr] (p2) -- (pool);
\draw[arr] (pool) -- (ret);
\end{tikzpicture}
\caption{Lifecycle of an item. Level~A (bridging) is only an upstream filter; the verdict comes
from the pilots (Level~B). A polarization rejection can be appealed directly to the pilots at the
cost of the author's reputation.}
\label{fig:pipeline}
\end{figure}
